\section{Conclusion and ongoing work}
\label{sec:conclusion}

We have described a real-world data-collection platform whose design is
governed by a single requirement: the recorded data must let a policy run
autonomously on the same rig. Two commitments follow from that requirement.
Every stream is timestamped on read and aligned to one reference time per
frame, with the residual drift logged, so temporal quality is measured
rather than assumed. And the recorded action is the constrained command the
robot actually executed, stored in both joint and end-effector form, with a
contract that inference replays a policy's output through the identical
command path.

The platform is measured and its defining constants travel with each
dataset, a readiness check confirms the rig is at a fixed collection state,
and the difficulties section keeps an honest record of what remains
imperfect --- chiefly the irreducible cross-camera skew and the several
thresholds still set by convention.

This is a living document. An experiments section will report the first
datasets collected on the platform, the measured drift budgets achieved,
and the autonomous performance of a joint-space and an end-effector policy
trained from the same episodes, once those runs exist.
